% !TEX program = xelatex
% ============================================================
%  خطة إعادة تصميم صفحة الإعدادات — تنظيم بصري + Accordion
%  الاستخدام على Overleaf:
%    1) New Project ← Blank Project
%    2) الصق المحتوى في main.tex
%    3) Menu ← Compiler ← اختر: XeLaTeX   (مهم للعربية)
%    4) Recompile ← نزّل PDF
% ============================================================
\documentclass[11pt]{article}

\usepackage{polyglossia}
\setmainlanguage{arabic}
\setotherlanguage{english}
\newfontfamily\arabicfont[Script=Arabic,Scale=1.0]{Amiri}
\newfontfamily\arabicfontsf[Script=Arabic,Scale=1.0]{Amiri}
\setmonofont{DejaVu Sans Mono}[Scale=0.8]

\usepackage{geometry}
\geometry{a4paper, margin=2.2cm}

\usepackage{xcolor}
\usepackage{graphicx}
\usepackage{booktabs}
\usepackage{longtable}
\usepackage{array}
\usepackage{enumitem}
\usepackage{titlesec}
\usepackage{fancyhdr}
\usepackage{hyperref}
\usepackage{tcolorbox}
\tcbuselibrary{skins,breakable}
\usepackage{fontawesome5}

% ── الألوان (هوية يَسِّر) ──
\definecolor{gold}{HTML}{B8933B}
\definecolor{darkbg}{HTML}{101414}
\definecolor{okgreen}{HTML}{27AE60}
\definecolor{warnred}{HTML}{C0392B}
\definecolor{softgray}{HTML}{F4F4F4}
\definecolor{midgray}{HTML}{7A7A7A}
\definecolor{codebg}{HTML}{1A2020}

\pagestyle{fancy}
\fancyhf{}
\renewcommand{\headrulewidth}{0.4pt}
\renewcommand{\footrulewidth}{0.4pt}
\fancyhead[R]{\small\color{gold}خطة تصميم صفحة الإعدادات — منصة يَسِّر}
\fancyhead[L]{\small\color{midgray}Settings UX Redesign}
\fancyfoot[C]{\small\thepage}

\titleformat{\section}{\Large\bfseries\color{gold}}{\thesection}{0.6em}{}
\titleformat{\subsection}{\large\bfseries\color{darkbg}}{\thesubsection}{0.5em}{}
\titleformat{\subsubsection}{\normalsize\bfseries\color{midgray}}{\thesubsubsection}{0.4em}{}

\setlist[itemize]{rightmargin=0pt, itemsep=2pt, topsep=3pt}
\setlist[enumerate]{rightmargin=0pt, itemsep=2pt, topsep=3pt}

\hypersetup{
  colorlinks=true, linkcolor=gold, urlcolor=gold,
  pdftitle={خطة تصميم صفحة الإعدادات}, pdfauthor={منصة يَسِّر}
}

% ── صناديق مخصّصة ──
\newtcolorbox{goalbox}[1][]{
  colback=gold!8, colframe=gold, boxrule=0.8pt,
  arc=3pt, left=8pt, right=8pt, top=6pt, bottom=6pt,
  breakable, #1
}
\newtcolorbox{notebox}[1][]{
  colback=softgray, colframe=midgray, boxrule=0.6pt,
  arc=3pt, left=8pt, right=8pt, top=6pt, bottom=6pt,
  breakable, #1
}
\newtcolorbox{warnbox}[1][]{
  colback=warnred!6, colframe=warnred, boxrule=0.8pt,
  arc=3pt, left=8pt, right=8pt, top=6pt, bottom=6pt,
  breakable, #1
}

\begin{document}

% ============================================================
\begin{center}
  {\Huge\bfseries\color{gold} خطة تحسين صفحة الإعدادات}\\[4pt]
  {\large\color{darkbg} تنظيم بصري + إعادة هيكلة بأقسام قابلة للطي}\\[6pt]
  {\small\color{midgray} منصة يَسِّر الفقهية — \today}
\end{center}

\vspace{4pt}
\hrule
\vspace{10pt}

% ============================================================
\section{الهدف}

\begin{goalbox}
تحويل صفحة الإعدادات من ست بطاقات متساوية متراصّة عموديًا (تمرير طويل وإجهاد بصري) إلى
\textbf{أربع مجموعات منطقية قابلة للطي (Accordion)} مع توحيد بصري شامل، وتقليل طول الصفحة،
وتسريع الوصول لكل إجراء.
\end{goalbox}

\subsection*{نطاق العمل}
\begin{itemize}
  \item \textbf{تحسين بصري:} ألوان، مسافات، ظلال، حركات، اتّساق الوضع الداكن.
  \item \textbf{إعادة هيكلة:} دمج وتجميع البطاقات في مجموعات Accordion منطقية.
  \item \textbf{بلا تغيير وظيفي:} منطق الأزرار والـ handlers يبقى كما هو بالكامل.
\end{itemize}

% ============================================================
\section{الوضع الحالي}

\begin{notebox}
\textbf{الملفان المتأثّران:}\\[3pt]
\texttt{src/Components/Views/SettingsView.tsx}\quad(854 سطرًا)\\
\texttt{src/Components/Views/SettingsView.module.css}\quad(700 سطر)\\
\texttt{src/Components/Views/SettingsView.test.tsx}\quad(قد يتأثر بالمُحدِّدات)
\end{notebox}

\vspace{6pt}
\textbf{البطاقات الست الحالية:}
\begin{enumerate}
  \item إصدار التطبيق والتحديثات
  \item تثبيت الموقع كتطبيق (PWA)
  \item المحتوى الفقهي وقاعدة المسائل
  \item إنجازاتي ومواظبتي
  \item الخصوصية والمزامنة السحابية
  \item الدعاء والتواصل مع المطوّر
\end{enumerate}

\textbf{المشاكل المرصودة:}
\begin{itemize}
  \item طول الصفحة يفرض تمريرًا طويلًا للوصول للأقسام السفلية.
  \item تكرار منطق شارات الحالة (\texttt{statusChip}) في 4 مواضع.
  \item تكرار كتل رسائل التغذية الراجعة (\texttt{feedbackAlert}) في 5 مواضع.
  \item استخدام إيموجي داخل دليل التثبيت يخالف اتّساق أيقونات الموقع.
\end{itemize}

% ============================================================
\section{المرحلة الأولى: إعادة الهيكلة}

تتحوّل البطاقات الست إلى \textbf{أربع مجموعات}، كل مجموعة Accordion مستقل:

\vspace{6pt}
\renewcommand{\arraystretch}{1.4}
\begin{longtable}{@{}p{0.05\textwidth} p{0.30\textwidth} p{0.48\textwidth} p{0.12\textwidth}@{}}
\toprule
\textbf{\#} & \textbf{المجموعة} & \textbf{تحتوي على} & \textbf{الافتراضي} \\
\midrule
\endhead
1 & التطبيق والتحديثات & إصدار التطبيق، تفريغ الكاش، تثبيت PWA & مفتوحة \\
2 & المحتوى الفقهي & المسائل والقاموس، مزامنة المحتوى & مفتوحة \\
3 & حسابي والخصوصية & إنجازاتي، المزامنة السحابية، مسح البيانات & مفتوحة \\
4 & الدعاء والتواصل & الدعاء + وسائل التواصل & مغلقة \\
\bottomrule
\end{longtable}
\renewcommand{\arraystretch}{1.0}

\begin{itemize}
  \item بطاقة \textbf{حساب الزائر} \texttt{GuestAccountCard} تبقى أعلى الصفحة خارج الـ Accordion.
  \item المجموعة 3 تعرض بطاقة تسجيل الدخول للزائر بدل خيارات المزامنة عند عدم وجود \texttt{user}.
  \item رأس الصفحة (الأيقونة + العنوان + الخط الفاصل) يبقى كما هو.
\end{itemize}

% ============================================================
\section{المرحلة الثانية: بنية الـ Accordion}

\textbf{مكوّن جديد \texttt{SettingsSection} داخل نفس الملف:}
\begin{itemize}
  \item \texttt{button} برأس قابل للنقر: أيقونة + عنوان + شارة حالة + سهم يدور $180^\circ$ عند الفتح.
  \item المحتوى بأسفله بـ انتقال ارتفاع سلس (\texttt{max-height} + \texttt{opacity}).
  \item دعم الإتاحة: \texttt{aria-expanded}, \texttt{aria-controls}, \texttt{role="region"}.
  \item حالة الفتح لكل مجموعة في \texttt{useState} منفصل، مع السماح بفتح أكثر من مجموعة.
  \item شارة صغيرة على رأس مجموعة «التطبيق» تظهر عند وجود تحديث متاح حتى وهي مطويّة.
\end{itemize}

\begin{goalbox}
\textbf{الفائدة:} تقليل طول الصفحة نحو $60\%$، وصول أسرع للأقسام، وتجربة أنظف على الجوال.
\end{goalbox}

% ============================================================
\section{المرحلة الثالثة: التحسين البصري}

\begin{enumerate}
  \item \textbf{بطاقات فرعية موحّدة:} خلفية أفتح وحدّ أنحف لإظهار التسلسل الهرمي داخل المجموعة.
  \item \textbf{رأس مجمّع:} خلفية متدرجة خفيفة وأيقونة بألوان الهوية عند الفتح (رمادي محايد عند الإغلاق).
  \item \textbf{توحيد \texttt{statusChip}:} مكوّن واحد بدل تكرار منطق الألوان في 4 مواضع.
  \item \textbf{توحيد \texttt{Feedback}:} مكوّن \texttt{<Feedback type text/>} واحد بدل 5 كتل متطابقة.
  \item \textbf{المسافات:} شبكة 8px منتظمة، وحدّ أقصى للعرض 720px للأقسام الطويلة.
  \item \textbf{الأيقونات:} استبدال إيموجي دليل التثبيت ($\Mobile$ $\Apple$ $\Desktop$) بأيقونات \texttt{react-icons}.
  \item \textbf{الوضع الداكن:} مراجعة الألوان الثابتة (\texttt{\#e74c3c}، \texttt{\#27ae60}) لتباعد كافٍ.
  \item \textbf{الحركة:} احترام \texttt{prefers-reduced-motion}، وتوحيد مدة الانتقالات على 250ms.
\end{enumerate}

% ============================================================
\section{المرحلة الرابعة: تحسينات تجربة المستخدم}

\begin{itemize}
  \item سهم المجموعة يتحرّك عند المرور ويشير إلى قابلية الفتح.
  \item الفتح الحر (لا طيّ تلقائي لما فوق المجموعة المفتوحة).
  \item زر «مسح البيانات» الأحمر يبقى داخل صندوق تأكيد لا يظهر إلا عند الطلب.
\end{itemize}

% ============================================================
\section{المرحلة الخامسة: الملفات والتحقق}

\textbf{ملفات تُعدَّل:}
\begin{itemize}
  \item \texttt{SettingsView.tsx} — إعادة هيكلة JSX + مكوّنا \texttt{SettingsSection} و \texttt{Feedback}.
  \item \texttt{SettingsView.module.css} — أنماط Accordion + التوحيد البصري.
\end{itemize}

\textbf{لا يُعدَّل:} منطق الأزرار والـ handlers (بلا تغيير وظيفي).

\vspace{6pt}
\textbf{أوامر التحقق:}
\begin{tcolorbox}[colback=codebg, colframe=codebg, coltext=white, arc=3pt, left=8pt, right=8pt, top=6pt, bottom=6pt]
{\ttfamily\small\raggedleft
npm run lint\\
npm run test -{}- SettingsView\\
npx tsc -{}-noEmit
}
\end{tcolorbox}

\begin{warnbox}
\textbf{المخاطر:} اختبار \texttt{SettingsView.test.tsx} يتوقّع عناصر بالنص؛ قد يحتاج تحديث المُحدِّدات
لتعكس البنية الجديدة دون كسر المنطق.
\end{warnbox}

% ============================================================
\section{معايير القبول}

\begin{itemize}
  \item[${\color{okgreen}\faCheckCircle}$] الصفحة تعرض أربع مجموعات Accordion قابلة للفتح والطي.
  \item[${\color{okgreen}\faCheckCircle}$] لا تغيير في سلوك أي زر أو إجراء بعد إعادة الهيكلة.
  \item[${\color{okgreen}\faCheckCircle}$] طول الصفحة المطويّة أقل بوضوح من الحالية.
  \item[${\color{okgreen}\faCheckCircle}$] الوضع الداكن سليم على كل العناصر الجديدة.
  \item[${\color{okgreen}\faCheckCircle}$] \texttt{lint} و \texttt{tsc} واختبارات \texttt{SettingsView} تمرّ بنجاح.
\end{itemize}

\end{document}
